\section{Introduction}
\label{sec:introduction}

A knowledge graph (KG) stores facts as triples. Each triple has a subject,
a relation, and a value. For example, a receipt graph can store a receipt
ID, the relation ``total'', and the value ``120.00''. Graphs built from
text often have missing values, wrong values, repeated triples, or invalid
relations. These errors affect search and later use of the data
\citep{ref_ji2021,ref_paulheim2017}. The source text gives us a way to check
and repair the graph.

Field rules and language models play two roles. Rules check the graph's
form. A model reads the text and proposes field values. KG quality tools
and SHACL provide rules for graph checks
\citep{ref_zaveri2016,ref_kontokostas2014,ref_shacl2017}.
Lin et al.\ add rules and graph context to LLM repair prompts
\citep{ref_lin2025}. We study three questions: when does repair improve a
graph, which steps produce the gain, and when does an existing graph help
more than building a new one?

We study graphs that store the fields of one document. Each graph has a
fixed document ID and a list of allowed relations. Each relation has at
most one value. Our method first cleans the input graph. A language model
then reads the source text and returns a complete graph. It can also read
an error report. A filter checks the document ID, allowed relations, text
matches, and field counts. It saves a reason for each rejected triple.

We compare three prompt settings: basic input, input with an error report,
and input with SHACL checks. They share the same instructions and output
limit. We score each response before and after filtering. This shows the
change from extra prompt text and the change from the filter. We also test
a selector that chooses edits one at a time and measure the cost of graph
checks.

The tests use three Chinese record collections, SROIE receipts, and CUAD
contracts. Chinese records list field names and values. Receipt text and
human field labels are stored apart. We use 20 new receipts to set the
source-line index, then fix the design and test it on 60 further receipts.
The contract test uses 56 documents. Two reviewers also check 200
reference differences and 200 system edits. These tests measure field
recovery, the loss of correct fields, and human acceptance of edits.

The paper makes three contributions:
\begin{itemize}
    \item A repair method that joins source-based generation with field
    checks and records why each triple is rejected.
    \item Tests on added defects, model outputs, receipts, and contracts,
    with scores for field recovery and retention and a human review of
    errors and edits.
    \item Shared-prompt and shared-response tests of error reports,
    source-line hints, SHACL context, and filtering. We also report full
    edit-selector runs and release the outputs and logs.
\end{itemize}

Section~\ref{sec:related_work} reviews earlier work.
Sections~\ref{sec:framework} and~\ref{sec:implementation} give the method
and its implementation. Section~\ref{sec:experiments} gives the tests and
results. Section~\ref{sec:conclusion} concludes.
